\section{Knowledge Base Construction and Governance}
\label{sec:knowledge_governance}

\subsection{Institutional Corpus and Source Provenance}
The factual foundation of KrishokChat is constructed from 2,946 verified institutional source documents published by accredited Bangladeshi agricultural research institutions:
\begin{itemize}
    \item \textbf{BARI Krishi Projukti Hatboi (9th Edition, 2023):} Comprehensive crop technology manuals covering field crops, tubers, and horticultural varieties \citep{bari_handbook2023}.
    \item \textbf{BRRI Adhunik Dhaner Chas (22nd Edition, 2024):} Canonical rice production and pathology guidelines \citep{brri_handbook2024}.
    \item \textbf{DAE Agrochemical Governance Registers (2024--2026):} Official lists of approved, restricted, and banned pesticides in Bangladesh.
\end{itemize}

\subsection{The 11-Slot Evidence Schema and Authority Stratification}
As detailed in Table~\ref{tab:knowledge_schema}, the knowledge base decouples semantic agronomic attributes (Layer A) from governance metadata (Layer B). To prevent conflicting recommendations, sources are stratified into a strict authority hierarchy:
\begin{equation}
    \text{MoA Regulatory} \succ \text{Current Institutional (BARI/BRRI)} \succ \text{Legacy Institutional} \succ \text{Secondary}
\end{equation}
When multiple sources contain overlapping entries, the resolution engine enforces precedence strictly by authority rank and publication timestamp, automatically invalidating superseded records.

\subsection{Cryptographic Hash-Chaining and Tamper-Evident Packs}
For offline and edge deployments, structured facts are compiled into an optimized SQLite database (the Offline Fact Pack). To guarantee data integrity, each record is hashed via SHA-256 and incorporated into a cryptographically signed Merkle manifest. During dynamic differential updates, client devices verify the cryptographic chain, ensuring \TamperDetRate{} detection of accidental corruption or adversarial tampering before committing updates.

\subsection{Knowledge Maintenance and Coverage Growth Loop}
To support continuous extension, KrishokChat implements a closed-loop knowledge maintenance workflow (Layer E24). When recurring queries trigger Tier 4 abstentions due to knowledge gaps, they are aggregated into an automated prioritization queue. Institutional agronomists author targeted tuples to bridge the gap. In empirical trials, an 18-minute targeted authoring session for Chili Anthracnose added 2 verified nodes, completely resolving 55 unanswerable queries (3.06 queries resolved per authoring minute) without requiring LLM re-training or full index rebuilding.

\begin{table*}[t]
\centering
\caption{Two-layer evidence schema: Agronomic Decision Contract (11 Semantic Slots) and Evidence Governance Layer.}
\label{tab:knowledge_schema}
\small
\begin{tabular}{llll}
\toprule
\textbf{Schema Layer} & \textbf{Field Name} & \textbf{Data Type / Format} & \textbf{Agronomic / Governance Purpose} \\
\midrule
\multirow{11}{*}{\textbf{Layer A: Semantic Contract}} 
 & \texttt{crop} & String (Canonical Bengali) & Target cultivated host crop (e.g., আলু / Potato) \\
 & \texttt{pathogen} & String (Canonical Bengali) & Target fungal, bacterial, or insect pest (e.g., নাবি ধসা / Late Blight) \\
 & \texttt{growth\_stage} & String (Categorical) & Permitted vegetative / reproductive phenological stage \\
 & \texttt{active\_ingredient} & String (Standard ISO) & Authorized active chemical compound (e.g., Mancozeb) \\
 & \texttt{formulation} & String (Standard Code) & Commercial formulation type (e.g., 80 WP, 250 EC) \\
 & \texttt{dose\_min} & Numeric (Float) & Minimum effective application rate per unit area/volume \\
 & \texttt{dose\_max} & Numeric (Float) & Maximum legally safe application rate without phytotoxicity \\
 & \texttt{dose\_unit} & String (SI Unit) & Measurement unit (e.g., g, ml, kg) \\
 & \texttt{solvent\_volume} & Numeric (SI Unit) & Dilution denominator (e.g., per 10 L water, per Decimal) \\
 & \texttt{interval\_days} & Integer (Days) & Minimum required interval between successive sprayings \\
 & \texttt{phi\_days} & Integer (Days) & Pre-Harvest Interval before crop consumption \\
\midrule
\multirow{6}{*}{\textbf{Layer B: Evidence Governance}}
 & \texttt{source\_authority} & Categorical (Level 1--4) & Regulatory (MoA) vs Institutional (BARI/BRRI) vs Secondary \\
 & \texttt{publication\_date} & ISO-8601 Date & Release date of the underlying institutional manual \\
 & \texttt{effective\_date} & ISO-8601 Date & Effective administrative enforcement date \\
 & \texttt{regulatory\_polarity} & Enum (\texttt{APPROVED}/\texttt{BANNED}) & National regulatory status under Bangladesh Pesticide Act \\
 & \texttt{provenance\_hash} & SHA-256 Hex & Cryptographic digest of source text and node metadata \\
 & \texttt{verifier\_signature} & Ed25519 Signature & Digital signature from institutional release authority \\
\bottomrule
\end{tabular}
\end{table*}

